\section{Beschreibung des bisherigen Workflows}
\label{ch:IST-Zustand:sec:Workflow}

Mithilfe der Ergebnisanzeige lässt sich der Allgemeinzustand des Systems überwachen (en. monitoring). Alle nachfolgenden Aussagen bezüglich des Workflows wurden aufgrund von Mitarbeiteraussagen (s. \autoref{Add:Mitschriften_Mitarbeiteraussagen}) sowie eigener Erfahrung im Einarbeitungsprozess zusammengefasst.

% 20 - 30 Tabs oder so - s. Quelle

Am Morgen öffnet die Person, die für das Monitoring zuständig ist, bis zu 20 Links für die verschiedenen Systeme und arbeitet diese stückweise ab.
Die Drop-Down"=Menüs der Attribute, die neben den Testergebnissen zu finden sind (s. vorheriger \autoref{ch:IST-Zustand:sec:Aufbau_altes_System}), werden zum Aufrufen der Oberflächen selten verwendet. Dadurch, dass alle vergangenen Varianten der Anwendungsfälle immer noch auswählbar und dadurch auch ungültige Kombinationen möglich sind, ist es gängig, sich alle benötigten Links als Lesezeichen in einem Ordner abzuspeichern.

Zum Monitoring wird jeweils nach gewissen Kennzahlen geschaut. Im Fehlerfall wird den zuständigen Personen Bescheid gegeben, damit sie diese die Probleme erforschen und schließlich beheben können.

Diese Vorgehensweise ist jedoch unpraktisch, da neue Mitarbeiter sich diese Links erst erfragen müssen, da diese nirgendwo einheitlich aufgelistet und eindeutig deklariert sind. Da meist nur der Link zur akut benötigten Seite erfragt wird, muss dieser vom Ansprechpartner entsprechend einzeln nach Anwendungsfall aus der Menge rausgefiltert werden, um diesen weiterzuleiten.

Weiterhin wird je nach Erfahrungsstand bzw. Kenntnissen entschieden, ob es sich am jeweiligen Tag um einen \textit{Green Day}\footnote{Ein \textit{Green Day} (dt. grüner Tag) ist ein Tag an dem das Testsystem keine Fehler anzeigt. Der Name kommt von der Farbdarstellung. Fehlerhafte Tests werden in der einzelnen Ergebnis-Übersicht rot dargestellt, während Tests, die fehlerfrei durchgelaufen sind, grün dargestellt werden.} handelt. In der nachfolgenden Nutzerstudie (s. \autoref{ch:Ergebnisse}) wurde von den Teilnehmern mehrfach angemerkt, dass es sich am entsprechenden Tag gemäß Er\-geb\-nis\-an\-zeige nicht um einen Green Day handle. Es würden jedoch aktuell nur bekannte Fehler der Testumgebung gezeigt werden, die nicht kurzfristig lösbar seien. Da diese nicht das Produkt selbst beträfen, werde es somit dennoch als Green Day gewertet.
